% Generated from locale/tr/content/many-valued-logic/sequent-calculus/rules-and-proofs.tex; do not hand-edit.


\olfileid[tr]{mvl}{seq}{rul}

\olsection{Kurallar ve \usetoken{P}{derivation}}

Aşağıda $\Gamma,\Delta,\Pi,\Lambda$, sonlu !!{sentence}s dizilerini
göstersin.

\begin{defn}[Sekant]
Bir \emph{$n$ taraflı sekant},
\[
\Gamma_1 \nSequent \dots \nSequent \Gamma_n
\]
biçiminde bir ifadedir; burada her $\Gamma_i$, $\Lang L$ dilinin
!!{sentence}s{}den oluşan sonlu, boş olması da mümkün, bir dizidir.
\end{defn}

\begin{defn}[Başlangıç sekantı]
Bir \emph{$n$ taraflı başlangıç sekantı}, dildeki herhangi bir !!{sentence}
$!A$ için $!A\nSequent\dots\nSequent !A$ biçiminde bir $n$ taraflı
sekanttır.

Dil $0$-konumlu bir $\star$ bağlacı, yani önerme sabiti içeriyorsa
$\star$'ın $\tf{\star}\in V$ ile ilişkili doğruluk değerinin konumunda
$\star$ bulunan ve diğer konumları boş olan
$\dots\nSequent\star\nSequent\dots$ sekantını da başlangıç sekantı sayarız.
\end{defn}

Bir $n$ değerli $\Log L$ mantığının her bağlacı için bu bağlacın $\Log L$
içinde alabileceği her doğruluk değerine karşılık gelen bir mantıksal kural
vardır. $\Log L$ için $n$ taraflı sekant hesabındaki !!^{derivation}s,
en üst sekantları başlangıç sekantları olan sekant ağaçlarıdır. Bir sekant
bir ya da daha çok sekantın altında duruyorsa $\Log L$ bağlaçlarına ait bir
çıkarım kuralıyla bunlardan doğru biçimde çıkmalıdır.

\begin{defn}[Teoremler]
!!^a{sentence} $!A$, $\Log L$'nin belirlenmiş doğruluk değerlerine karşılık
gelen her konumda $!A$ bulunan $n$-sekantın !!a{derivation} varsa, $n$
değerli $\Log L$ mantığının bir \emph{teoremidir}. $!A$ teoremse
$\Proves[\Log L]!A$, değilse $\Proves/[\Log L]!A$ yazarız.
\end{defn}

\begin{defn}[!!^{derivability}]
Bir !!^{sentence} $!A$, $\Log L$ $n$ değerli mantığında bir !!{sentence}s
kümesi $\Gamma$'dan \emph{!!{derivable}} ise ve ancak o zaman
$\Gamma\Proves[\Log L]!A$ yazarız. Bunun anlamı şudur: sonlu bir
$\Gamma_0\subseteq\Gamma$ ve $\Gamma_0$ içindeki !!{sentence}s{}in oluşturduğu
bir $\Gamma_0'$ dizisi vardır ve
\[ \Lambda_1 \nSequent \dots \nSequent \Lambda_n \]
sekantının !!a{derivation} vardır; burada $\Lambda_i$, $i$'inci konum
belirlenmiş bir doğruluk değerine karşılık geliyorsa $!A$, aksi durumda
$\Gamma_0'$'dür. $!A$, $\Gamma$'dan !!{derivable} değilse
$\Gamma\Proves/!A$ yazarız.
\end{defn}

Örneğin üç değerli Łukasiewicz mantığının üç taraflı bir sekant hesabı
vardır. $\Gamma\nSequent\Pi\nSequent\Delta$ sekantında $\Gamma$, $\False$;
$\Delta$, $\True$; $\Pi$ ise $\Undef$ değerine karşılık gelir. Aksiyomlar
$!A\nSequent !A\nSequent !A$ biçimindedir. Yalnızca $\True$ belirlenmiş
olduğundan $\Gamma\Proves[\LogLuk[3]]!A$ ancak ve ancak
$\Gamma\nSequent\Gamma\nSequent !A$ sekantının !!a{derivation} varsa
geçerlidir. $\Undef$ da belirlenmiş olsaydı
$\Gamma\nSequent !A\nSequent !A$ sekantının !!a{derivation} gerekirdi.

